% !TeX root = ../教材题库_全部条目_转写工作稿.tex
% Source: XJ-B1-C06-EX-P237-Y1511
% 第6章 / 6.4 用样本估计总体 / 教材 p.237、238 / PDF 第244、245页
% OCR未校对：公式与文字均可能有误；未完成数学TeX转写；禁止直接发给学生
% 原题及相邻原解截图（核对用）：../题目原图/XJ-B1-C06-EX-P237-Y1511-01.webp, ../题目原图/XJ-B1-C06-EX-P237-Y1511-02.webp
% 本文件只在转写完成后作为题干使用；OCR 条目尚不是可靠题干。
\begin{bankitem}{6.4 例题 例题3}{第6章 / 6.4 用样本估计总体 / 教材 p.237、238 / PDF 第244、245页}{XJ-B1-C06-EX-P237-Y1511}{OCR未校对：公式与文字均可能有误；未完成数学TeX转写；禁止直接发给学生}
\begin{ocrdraft}
[教材p.237 / PDF第244页]
\par
0河某市进行家庭人均年收人调查时，分别对城镇家庭和农村家庭进行调
\par
查．在全部城镇的285679户中无放回地随机抽取了350户，在全部农村的275692
\par
户中无放回地随机抽取了350户．调查结果为：城镇家庭人均年收人是54188元，
\par
[教材p.238 / PDF第245页]
\par
农村家庭人均年收人是23119元．试估计该市家庭人均年收人．
\par
解此抽样为分层抽样，但不是同比例抽样．设总体，A表示该市所有家庭的人
\par
均年收人，总体A分为两层：第一层对应所有城镇家庭的人均年收人，第二层
\par
，A2对应所有农村家庭的人均年收人．
\par
用五表示来自总体A的样本均值，用五表示来自总体，A2的样本均值，则一
\par
54188，2=23119.
\par
285679
\par
A在A中所占的比例是WI一
\par
0．5089．
\par
285679十275692
\par
275692
\par
，A2在A中所占的比例是W2一
\par
0．4911.
\par
285679十275692
\par
所以A的总体均值的估计是
\par
X=VVlxl十W2X2==0.5089\ensuremath{\times}54188十0．4911\ensuremath{\times}2311938930．
\par
即该市家庭人均年收人的估计是38930元．
\par
我们知道，在分层抽样中，用N表示总体，A的个体总数，若将总体A分为L层
\par
(L<N),用N，表示第》层0一1，2，
\par
，L）的个体总数，则有N=NI十N2十
\par
，十
\par
NI\_．记、铲一一0一1，2，
\par
，L)，称WI为第》层的权重．
\par
对》一1，2，
\par
，L，用表示从第层抽出样本的均值．我们称
\par
X==VVlxl十、V2工2十．，十WI\_工L
\par
是总体均值的估计．
\par
分层抽样在获得总体均值估计的同时，也得到各层的均值估计．在例3中，不
\par
但得到了A的均值估计，还得到了AI和，A2的均值估计．
\end{ocrdraft}
\end{bankitem}
